\subsection{Round two: bitwise seeds and structure analyses}
\label{sec:r2-method}

Round one seeded VPR with module-level clusters. Datapath logic is bit-level: a
32-bit register or adder is a regular stack of identical bit cells, and the carry chain
already shows that VPR places such stacks as ordered vertical runs. Round two therefore
adds a bit-level seeding aspect that keeps the module hierarchy at the coarse level
(Section~\ref{sec:r2-appendix} records where the idea came from).

\paragraph{Word membership from Odin~III data, not names.}
The obvious source, the provenance BLIF's block name \code{module/\$nN@file:line}, is too
coarse: the line is the HDL \emph{statement}, so one \code{always} block is one group (the
largest (module, file:line) group in any design holds up to 1739 blocks,
and sha's holds 1739 of its 2506 blocks; 59 of the 161 statement lines that
hold word blocks hold more than one word, up to 51). The finer records exist in the Odin~III session that
writes the BLIF, so the seed builder reruns the identical script, checks that the BLIF it
writes equals the kept one byte for byte (all 9 designs
match), and reads the records Loki~III's labeller reads: the HDL name \code{bus[i]} of a
cell's output net and the source spans of each cell, matched to the word-level cells that
exist before technology mapping. A final cell joins its BLIF block through the net id the BLIF
name carries (\code{\$n<id>}), so the join is by identifier, not by string
(241 of 11198 cells, 2.2\%, found no BLIF block and are left unlabelled;
the writer legalises dotted names, handled explicitly). A block's word is the first rule
that applies: (bus) an output net named \code{bus[i]} puts it in word \code{bus} at bit $i$,
which holds for every flip-flop that keeps a register name; (cell) exactly one word-level cell's
spans lie inside the block's spans, which is true for adders, multipliers, memories and
single-statement LUTs; otherwise (none) the block has no word. Bit order inside an adder word is
the carry-chain position. ABC's cone provenance merges the spans of a LUT that mixes words,
and its edge budget makes some runs share one record, so 21\% of the LUTs get a word
against 100\% of the flip-flops and 99\% of the hard cells. Over all designs,
58.7\% of the blocks are in a word (per-design median
63.1\%); 444 words in all.

\paragraph{Seed families.}
A CLB joins a word cluster when at least one third of its atoms carry a word and that word is
the plurality; otherwise it keeps its Urd module cluster. Four families follow, each entering
VPR as a soft initial placement at VPR's own start temperature (scale~1).
\code{words}: the word CLB groups are the clusters, arranged by the round-one quotient-graph
layout. \code{slices}: the same clusters, but each word is laid out as a vertical run in bit
order (bit~0 on top, the carry-chain direction), runs of equal length aligned to the same rows,
so that bit $i$ of connected words shares a row; true carry chains are untouched.
\code{hier}: Urd module clusters at the coarse level and, inside each module, the words as
bit-sliced runs. \code{words\_partial}: only the word CLBs are seeded, every other CLB keeps
VPR's analytic position (and every carry chain with an unseeded CLB stays whole). The
\code{words} family also has the inter-cluster-only variant. Pads and multipliers stay
anchored as in round one.

\paragraph{Random and coarse-random arms (RQ-N4).}
\code{random} is a random legal start: every logic CLB takes a random slot and carry chains a
random window, fed as a soft seed at scale~1. The coarse-random arm keeps the clusters and
their within-cluster order but permutes the cluster positions on the die (\code{urd\_shuf} for
Urd, and the \code{*\_shuf} twin of the best new family).

\paragraph{Structure analyses (no VPR runs).}
RQ-N1 measures, on each design's unseeded baseline, how module- and word-coherent three
placements are: VPR's analytic start, its final placement, and a random null (logic-CLB
coordinates permuted, 20 samples). Coherence is the NMI between a $g\times g$ grid partition
($g=\lceil\sqrt{\#\text{labels}}\rceil$) and the labels, and the compactness
(label CLBs per grid tile of the label's bounding box, CLB-weighted). RQ-N2 correlates the seed
gain with three structure measures of the netlist graph: Louvain modularity $Q$; the $z$-score of
$Q$ against a degree-preserving null (double edge swaps, 20 samples); and the inverse Rent
exponent $-p$, with $p$ fitted from the terminals of BFS balls of geometric sizes (a lower
exponent is a more local design). Spearman's $\rho$ and $p$-value are reported with $n=9$,
which is small.

\paragraph{Hypotheses and decision rules.}
Paired comparisons are per design (variant over reference), reported as a geomean ratio, the number of
designs where the variant is better and a two-sided exact sign test. H7a: \code{words} beats Urd in
soft wirelength (supported if the geomean is below 1 and $p<0.05$, not supported if the
geomean is at least 1, else inconclusive). H7b: \code{slices} beats \code{words} on wirelength and
channel width (supported if both geomeans are below 1 and the wirelength sign test has $p<0.05$).
H7c: \code{hier} has the lowest geomean wirelength of all eleven soft families. H7d: \code{words\_partial}
is at least as good as the full word seed (supported if the geomean ratio is at most 1, inconclusive up to 1.01).
H1: the baseline's final placement is module- and word-coherent (supported if, for both labellings,
at least two thirds of the designs lie more than two null standard deviations above the null in both
NMI and compactness). H2: structure predicts gain (supported if at least two of the three
measures correlate positively with the best round-two gain at $p<0.05$; not supported if no
correlation reaches $p<0.05$). H4: random start $>$ permuted clusters $>$ quotient layout $\ge$ analytic start in
wirelength ratio, for Urd and for the best new family.
